% !PRINTED none
% Abschnitt-title page (part divider): no running head/page number visible in
% the scan for this page; §. 169 opens on it.

\begin{center}
\gesperrt{Twenty-first Chapter.}

\gesperrt{Cyclotomic Fields.}
\end{center}

\begin{center}
\rule{4cm}{0.4pt}
\end{center}

\begin{center}
\S.\,169.

\gesperrt{Decomposition of the prime $q$ into prime factors in the cyclotomic field $\Omega_{q^{\varkappa}}$.}
\end{center}

We make a second application of the general theory of algebraic numbers to the cyclotomic fields, and thereby gain the means of bringing the theory of the \gesperrt{Abelian} number fields in general to a beautiful conclusion.

The considerations that we have first to undertake can be applied, with small modifications, to every full cyclotomic field (\S.\,18). For the sake of simplicity, however, we restrict ourselves here to the case in which the degree of the root of unity is a power of a prime, which suffices for the intended application.

Let $q$ be a natural prime (with the inclusion of 2) and
\weq{1}{m = q^{\varkappa}}
a power of $q$, whose positive exponent $\varkappa$ is assumed to be greater than 1 for $q=2$.

Let $r$ further be a primitive $m$th root of unity and $\Omega_m$ the full cyclotomic field for the exponent $m$, whose degree is
\weq{2}{\mu = \varphi(m) = q^{\varkappa-1}(q-1)}
(Bd. I, \S.\,134).

We denote throughout by $n$ the $\mu$ numbers of a complete system of residues for the modulus $m$, with the exclusion of the numbers divisible by $q$; and the $\mu$ numbers $r^n$ are then the roots of the irreducible equation $f(x)=0$ of degree $\mu$, in which
\weq{3}{f(x) = \dfrac{x^m-1}{x^{\frac{m}{q}}-1} = x^{q^{\varkappa-1}(q-1)} + x^{q^{\varkappa-1}(q-2)} + \cdots + 1}
%CONT>>
